\pdfminorversion=7%
\documentclass[aspectratio=169,mathserif,notheorems]{beamer}%
%
\xdef\bookbaseDir{../../bookbase}%
\xdef\sharedDir{../../shared}%
\RequirePackage{\bookbaseDir/styles/slides}%
\RequirePackage{\sharedDir/styles/styles}%
\toggleToGerman%
%
\subtitle{19.~Fabrik-Datenbank:~Tabellen und Sichten in LibreOffice Base}%
%
\begin{document}%
%
\startPresentation%
%
\section{Einleitung}%
%
\begin{frame}%
\frametitle{Datenbankzugriff mit Python}%
%
\begin{itemize}%
\item Bisher können wir unsere \glslink{db}{Datenbanken} nur mit \glsShort{SQL} bearbeiten.%
%
\item<2-> Das ist ungeignet für die große Mehrheit der Leute in einer Organisation.%
%
\item<3-> Leute sind daran gewöhnt, mit \microsoftExcel-Tabellen oder \microsoftWord-Dokumenten zu arbeiten.%
%
\item<4-> Wenn wir ihnen sagen, dass sie eine Programmiersprache für Daten lernen sollen, wird das auf wenig Gegenliebe stoßen.%
%
\item<5-> Aber nun werden wir eine andere Möglichkeit, mit unserer \glslink{db}{Datenbank} zu interagieren lernen.%
\end{itemize}%
\end{frame}%
%
%
\section{Tabellen in LibreOffice Base}%
%
\begin{frame}[t]%
\frametitle{Daten über LibreOffice Base in Tabellen eingeben}%
%
\begin{itemize}%
%
\only<-2>{%
\item Wir wollen \libreofficeBase\ verwenden, um mit einer Tabelle in unserer \postgresql\ \glslink{db}{Datenbank} zu arbeiten.%
}%
%
\only<-2>{%
\item<2-> Wir öffnen also unser Dokument \textil{factory.odb} aus der vorigen Einheit mit \libreofficeBase\ und verbinden uns auf unsere \glslink{db}{Datenbank}, wobei wir das Passwort \textil{superboss123} eingeben müssen.%
}%
%
\only<-3>{%
\item<3-> Wir nun wählen die Tabelle \textil{demand} in der Tabellen-Sicht aus und Doppelklicken darauf.%
}%
%
\only<-4>{%
\item<4-> Ein neues Fenster öffnet sich.%
}%
%
\only<-5>{%
\item<5-> In diesem Fenster sehen wir die Daten unserer Tabelle.%
}%
%
\only<-6>{%
\item<6-> Die Spaltennamen unserer Tabelle sind auch hier Spaltennamen.%
}%
%
\only<-7>{%
\item<7-> Jeder Datensatz ist auf einer Zeile dargestellt.%
}%
%
\only<-8>{%
\item<8-> Das ist schon mal sehr gut.%
}%
%
\only<-9>{%
\item<9-> Wie gesagt, viele Leute sind mit \microsoftExcel-Tabellen vertraut.%
}%
%
\only<-10>{%
\item<10-> Die Ansicht hier sieht dem schon recht ähnlich.%
}%
%
\only<-11>{%
\item<11-> Die Benutzer können die Konzepte von Spalten und Zeilen intuitiv verstehen.%
}%
%
\only<-12>{%
\item<12-> Das ist viel klarer, also die Ausgaben von \psql\dots%
}%
%
\only<-13>{%
\item<13-> Wir können die Daten hier auch gleich bearbeiten.%
}%
%
\only<-14>{%
\item<14-> Wir können aber auch neue Datensätze eingeben.%
}%
%
\only<-15>{%
\item<15-> Dafür positionieren wir den Kursor in das zweite Feld der leeren Zeile am Ende der Tabelle.%
}%
%
\only<-16>{%
\item<16-> Wir überspringen die Spalte \sqlil{id}, weil ihr Wert automatisch vom \glsShort{dbms} gesetzt wird.%
}%
%
\only<-17>{%
\item<17-> Wir wählen Mr.~Bobbo aus, geben also customer ID~4 ein.%
}%
%
\only<-18>{%
\item<18-> Wir speichern das er am 12.~April 2025 11 Einheiten des Produkts mit ID~7, also Schuhe der Größe~37, bestellt hat.%
}%
%
\only<-19>{%
\item<19-> Nachdem wir diese Daten eingegeben haben und unser Kursor am Ende der letzten Spalte ist, drücken wir \keys{\tab}.%
}%
%
\only<-20>{%
\item<20-> Die Daten werden nun an das \glsShort{dbms} geschickt.%
}%
%
\only<-21>{%
\item<21-> Sehen Sie, dass im Feld \sqlil{id} des neuen Datensatzes eine~0 steht?%
}%
%
\only<-22>{%
\item<22-> Der Grund ist, dass wir hier keinen Wert eingegeben haben.%
}%
%
\only<-23>{%
\item<23-> Das System hat den neuen Datensatz an das \glsShort{dbms} geschickt.%
}%
%
\only<-24>{%
\item<24-> Das \glsShort{dbms} setzt das \sqlil{id}-Feld automatisch.%
}%
%
\only<-25>{%
\item<25-> Die \libreofficeBase\ \glsShort{GUI} weis das aber nicht\dots%
}%
%
\only<-26>{%
\item<26-> Um den wirklichen Wert des Feldes~\sqlil{id} zu sehen, müssen wir die Daten nochmal neu laden.%
}%
%
\only<-27>{%
\item<27-> Wir klicken auf den Neu-Laden-Button~\libreOfficeBaseRefresh.%
}%
%
\only<-28>{%
\item<28-> Danach hat das \sqlil{id}-Feld unseres neuen Datensatzes den korrekten Wert~13.%
}%
%
\item<29-> Wir schließen das Tabellen-Fenster und gehen zurück ins Hauptfenster.%
%
\end{itemize}%
%
\locateGraphicTB{3}{width=0.54\paperwidth}{graphics/tableAndView/factoryLibreOfficeBaseTableAndView1chooseTableDemand}{0.23}{0.22}%
\locateGraphicTB{4-15}{width=0.54\paperwidth}{graphics/tableAndView/factoryLibreOfficeBaseTableAndView2tableDemand}{0.23}{0.22}%
\locateGraphicTB{16-17,20}{width=0.54\paperwidth}{graphics/tableAndView/factoryLibreOfficeBaseTableAndView3addDemand}{0.23}{0.22}%
\locateGraphicTB{18-19}{width=0.54\paperwidth}{graphics/tableAndView/factoryLibreOfficeBaseTableAndView3addDemand}{0.4}{0.22}%
\locateGraphicTB{21-27}{width=0.54\paperwidth}{graphics/tableAndView/factoryLibreOfficeBaseTableAndView4demandAdded}{0.23}{0.22}%
\locateGraphicTB{28-}{width=0.54\paperwidth}{graphics/tableAndView/factoryLibreOfficeBaseTableAndView5demandsRefreshed}{0.23}{0.22}%
%
\end{frame}%
%
\begin{frame}
\frametitle{Bestellungen vorher}%
\gitLoadAndExecSQL{factory:select_sale_data_1}{}{factory}{select_sale_data.sql}{factory}{boss}{superboss123}%
\listingSQLandOutput{}{factory:select_sale_data_1}{}{0.05}{0.11}{0.9}{0.89}%
\end{frame}
%
\begin{frame}
\frametitle{Bestellungen nachher}%
\gitLoadAndExecSQL{factory:update_according_to_base_table_add_tab12}{}{factory}{update_according_to_base_table_add.sql}{factory}{boss}{superboss123}%
\gitLoadAndExecSQL{factory:select_sale_data_2}{}{factory}{select_sale_data.sql}{factory}{boss}{superboss123}%
\listingSQLandOutput{}{factory:select_sale_data_2}{}{0.05}{0.11}{0.9}{0.89}%
\end{frame}
%
%
\section{Sichten in LibreOffice Base ausführen}%
%
%
\begin{frame}[t]%
\frametitle{Sichten in LibreOffice Base ausführen}%
%
\begin{itemize}%
%
\only<-3>{%
\item Eine Datenbanksicht stellt einer Applikation eine Perspektive auf die Daten zur Verfügung.%
%
\item<2-> Eine Sicht sieht ein bisschen aus wie eine Tabelle, ist aber eher sowas wie eine gespeicherte \glsShort{SQL}-Abfrage.%
%
\item<3-> Wir können auch aus \libreofficeBase\ auf Sichten zugreifen {\dots} und dort sehen sie auch wie Tabellen aus.%
}%
%
\only<-5>{%
\item<4-> Zurück in der \inQuotes{Tables} Ansicht klicken wir doppelt auf die Siche \sqlil{sale}.%
}%
%
\only<-6>{%
\item<5-> Dadurch öffnet sich wieder ein neues Fenster.%
}%
%
\only<-7>{%
\item<6-> Wie sehen jetzt das Ergebnis der Abfrage in einer schönen tabellarischen Form.%
}%
%
\only<-7>{%
\item<7-> Vielleicht ist die Spalte \sqlil{customer_name} dabei ein Wenig komisch auf Ihrem Rechner dargestellt.%
}%
%
\only<-9>{%
\item<8-> In dem Fall können Sie ihre Größe ändern und etwas breiter machen, in dem Sie ihren rechten Rand mit der linken Maustaste ziehen.%
}%
%
\only<-10>{%
\item<9-> Dann erscheinen alle Daten korrekt.%
}%
%
\only<-11>{%
\item<10-> Vorhin hatten wir ja neue Daten in die Tabelle für Bestellungen eingetragen.%
}%
%
\only<-12>{%
\item<11-> Früher hatten wir nur eine Bestellung für Mr.~Bobbo in unserem System.%
}%
%
\only<-13>{%
\item<12-> Nun erscheint eine Neue in der letzten Zeile der Sicht.%
}%
%
\item<13-> Das bestätigt noch einmal, dass unsere Sicht neu ausgeführt wurde und das unsere Änderungen tatsächlich in der Datenbank gespeichert wurden.%
%
\end{itemize}%
%
\locateGraphicTB{4}{width=0.54\paperwidth}{graphics/tableAndView/factoryLibreOfficeBaseTableAndView6chooseViewSale}{0.23}{0.3}%
\locateGraphicTB{5-}{width=0.54\paperwidth}{graphics/tableAndView/factoryLibreOfficeBaseTableAndView7viewSale}{0.23}{0.3}%
%
\end{frame}%
%
\section{Zusammenfassung}%
%
\begin{frame}%
\frametitle{Zusammenfassung}%
\begin{itemize}%
\item An dieser Stelle sehen wir bereits einen wichtigen Vorteil davon, eine \glsShort{GUI} wie \libreofficeBase\ zu verwenden.%
%
\item<2-> Wir können nun die Daten in unseren Tabellen viel bequemer eingeben und anzeigen.%
%
\item<3-> Vorher mussten wir uns die abstrakte tabellarische Form unserer Tabellen mit unserer Fantasie vorstellen.%
%
\item<4-> Aber jetzt können wir sie wirklich als Tabellen sehen und editieren.%
%
\item<5-> Dadurch haben wir ein viel besseres Look \& Feel für unsere Datenbank.%
%
\item<6-> Mit so einem Interface könnte ein normaler Benutzer bereits arbeiten.%
%
\item<7-> Das ist schon sehr nett.%
\end{itemize}%
\end{frame}%
%
\endPresentation%
\end{document}%%
\endinput%
%
